%%%%%%%%%%%%%%%%%%%%%%%%%%%%%%%%%%%%%%%%%%%%%%%%%%%%%%%
%%         BORRAR EL CONTENIDO DEL DOCUMENTO         %%
%%   SOLO ILUSTRA CONTENIDO QUE EJEMPLIFICA EL USO   %%
%%%%%%%%%%%%%%%%%%%%%%%%%%%%%%%%%%%%%%%%%%%%%%%%%%%%%%%

%\subsection{Abtsract}
\lipsum[1]

\section{La sección}
\lipsum[1]

\subsection{La subsección}
\lipsum[1]

\subsubsection{La subsubsección}
\lipsum[1]

%\RaggedRight
\justifying
\section*{Elementos comúnes}
\markboth{Elementos comúnes}{} % Agregar cuando no se enumera la sección

\subsection{Espacio de código}
\begin{lstlisting}[language=Java]
    /**
     * Habilita una categoría
     * @param id identificador de la categoría (visto en el path)
     * @return Una respuesta de éxito o fallo
     * @mapeo /category/{id}/enable
     * @metodo HTTP PATCH
     * @estado 200 - La categoría ha sido activada
     * @estado 400 - Ocurrió un error inesperado
     * @estado 404 - ID no encontrado
     */
    @Operation(summary = "Habilita una categoría", description = "Cambia el estado de una categoría a activo (1)")
    @ApiResponses(value = {
        @io.swagger.v3.oas.annotations.responses.ApiResponse(
            responseCode = "200",
            description = "La categoría ha sido activada",
            content = @Content(
                mediaType = "text/plain",
                schema = @Schema(type = "string"),
                examples = @ExampleObject(value = "La categoría ha sido activada")
            )
        ),
        @io.swagger.v3.oas.annotations.responses.ApiResponse(
            responseCode = "400",
            description = "Ocurrió un error inesperado",
            content = @Content(
                mediaType = "text/plain",
                schema = @Schema(type = "string")
            )
        ),
        @io.swagger.v3.oas.annotations.responses.ApiResponse(
            responseCode = "404",
            description = "ID no encontrado",
            content = @Content(
                mediaType = "text/plain",
                schema = @Schema(type = "string"),
                examples = @ExampleObject(value = "El id de la categoría no existe")
            )
        )
    })
    @PatchMapping("/{id}/enable")
    public ResponseEntity<ApiResponse> enable(@PathVariable Integer id) {
        return ResponseEntity.ok(svc.enable(id));
    }
\end{lstlisting}
%
\subsection{Tabla}
\figura[H]{
    \begin{tabla}![Variante4][Variante5]{| m{0.3\linewidth} | m{0.3\linewidth} | m{0.3\linewidth} |}
        \lipsum[1][1] & \lipsum[1][1] & \lipsum[1][1] \\ \hline
        \lipsum[1][1] & \lipsum[1][1] & \lipsum[1][1] \\ \hline
        \lipsum[1][1] & \lipsum[1][1] & \lipsum[1][1] \\ \hline
        \lipsum[1][1] & \lipsum[1][1] & \lipsum[1][1] \\ \hline
        \lipsum[1][1] & \lipsum[1][1] & \lipsum[1][1] \\ \hline
        \lipsum[1][1] & \lipsum[1][1] & \lipsum[1][1] \\ \hline
    \end{tabla}
}{Tabla Muestra}{fig:tabla}

\subsection{Imagen}\label{sec:imagen}
La imagen \Fig{fig:gato} de la \Sec{sec:imagen}, en efecto es un gato.
\fimage[0.5]{cat}{Gato}{fig:gato}
%
\section{Listas}
\subsection{Itemize}
\begin{itemize} [label = \itz]
    \item Categoría 1
    \begin{itemize} [label = \iitz]
        \item Subcategoría 1
            \begin{itemize} [label = \iiitz]
                \item Subsubcategoría 1
            \end{itemize}
        \item Subcategoría 2
    \end{itemize}
    \item Categoría 32
\end{itemize}
    
\subsection{Enumerate}
\begin{enumerate} [label = \ienum]
    \item Categoría 1
    \begin{enumerate} [label = \iienum]
        \item Subcategoría 1
            \begin{enumerate} [label = \iiienum]
                \item Subsubcategoría 1
            \end{enumerate}
        \item Subcategoría 2
    \end{enumerate}
    \item Categoría 32
\end{enumerate}
%
\section{Extras}
\subsection{Operaciones con align}
\begin{align*}
    \frac{n}{2^1} + \frac{n}{2^2} + ... + \frac{n}{2^p} 
    &= (\frac{n}{2^0} + \frac{n}{2^1} + \frac{n}{2^2} + ... + \frac{n}{2^p}) - \frac{n}{2^0}\\
    &= (\frac{n}{1} + \frac{n}{2^1} + \frac{n}{2^2} + ... + \frac{n}{2^p}) -\frac{n}{1}\\
    &= (n + \frac{n}{2^1} + \frac{n}{2^2} + ... + \frac{n}{2^p}) - n\\
    &= (\sum_{i=0}^{p} \frac{n}{2^i}) - n \\
    &= n(\sum_{i=0}^{p} \frac{1}{2^i}) - n \\
    &= n(2 - 2^{-p}) - n\\
    &= n(2 - \frac{1}{2^{p}}) - n\\
    &= n(2 - \frac{1}{n})\\
    &= 2n - \frac{n}{n} - n\\
    &= 2n - 1 - n\\
    &= n - 1
\end{align*}

\subsection{Tableros, para ilustraciones o figuras}
{
\centering
\subsubsection{Ajedrez}
\renewcommand{\arraystretch}{2}
\setlength{\arrayrulewidth}{1.5pt}
\begin{NiceTabular}{| >{\centering\arraybackslash} m{0.65cm} 
                      >{\centering\arraybackslash} m{0.65cm} 
                      >{\centering\arraybackslash} m{0.65cm} 
                      >{\centering\arraybackslash} m{0.65cm} 
                      >{\centering\arraybackslash} m{0.65cm} 
                      >{\centering\arraybackslash} m{0.65cm} |
                    }
    \CodeBefore
        \chessboardcolors{Foreground}{Background}
    \Body
        \hline
        & & & & & \\
        & & & & & \\
        & & & & & \\
        & & & & & \\
        & & & & & \\
        & & & & & \\ \hline
\end{NiceTabular}

\subsubsection{Dominó}
\begin{NiceTabular}{| >{\centering\arraybackslash} m{0.65cm} 
                      >{\centering\arraybackslash} m{0.65cm} 
                    | >{\centering\arraybackslash} m{0.65cm} 
                      >{\centering\arraybackslash} m{0.65cm} 
                    | >{\centering\arraybackslash} m{0.65cm} 
                      >{\centering\arraybackslash} m{0.65cm} |
                    }
    \CodeBefore
        \rowcolor{Background}{1,2,3,4,5,6}
        \cellcolor{Variante1}{1-1,1-2,1-5,1-6,2-3,2-4,3-1,3-2,3-5,3-6,4-3,4-4,5-1,5-2,5-5,5-6,6-3,6-4}
    \Body
        \hline
        & & & & & \\
        & & & & & \\ \hline
        & & & & & \\
        & & & & & \\ \hline
        & & & & & \\
        & & & & & \\ \hline
\end{NiceTabular}

\subsubsection{Programación dinámica}
\begin{NiceTabular}{| >{\centering\arraybackslash} m{0.5cm} 
                    | >{\centering\arraybackslash} m{0.5cm} 
                    | >{\centering\arraybackslash} m{0.5cm} 
                    | >{\centering\arraybackslash} m{0.5cm} 
                    | >{\centering\arraybackslash} m{0.5cm} |
                    }[first-row,first-col]
    \CodeBefore
        \chessboardcolors {Variante1}{Background}
        \columncolor{Variante1!25!Background}{0}
        \rowcolor{Variante1!25!Background}{0}
        \cellcolor{Variante6!45!Background}{3-3}
        \cellcolor{Background}{0-1}
        \cellcolor{Background}{1-0}
    \Body
          &   & J & O & J & I \\ \hline
          & 0 & 0 & 0 & 0 & 0 \\ \hline
        R & 0 & 0 & 0 & 0 & 0 \\ \hline
        O & 0 & 0 & 1 & 0 & 0 \\ \hline
        J & 0 & 1 & 0 & 2 & 0 \\ \hline
        O & 0 & 0 & 1 & 0 & 0 \\ \hline
\end{NiceTabular}
    
\renewcommand{\arraystretch}{1}
\setlength{\arrayrulewidth}{1pt}
}